% Generated draft: requires mathematical and source review.
\documentclass[11pt]{article}
\usepackage[margin=25mm]{geometry}
\usepackage{fontspec}
\setmainfont{Latin Modern Roman}
\usepackage{amsmath,amssymb,mathtools,unicode-math}
\setmathfont{Latin Modern Math}
\usepackage{enumitem,needspace,xurl,hyperref}
\setlength{\parindent}{0pt}
\setlength{\parskip}{5pt}
\setlength{\emergencystretch}{3em}
\newcommand{\N}{\mathbb{N}}
\newcommand{\Z}{\mathbb{Z}}
\newcommand{\Q}{\mathbb{Q}}
\newcommand{\R}{\mathbb{R}}
\newcommand{\C}{\mathbb{C}}
\newcommand{\F}{\mathbb{F}}
\newcommand{\A}{\mathbb{A}}
\newcommand{\PP}{\mathbb{P}}
\newcommand{\E}{\mathbb{E}}
\newcommand{\eps}{\varepsilon}
\newcommand{\dd}{\,\mathrm{d}}
\begin{document}
\section*{UnsolvedMath ID: 701698}
\textbf{Source problem code:} catalogue entry
\subsection{Definitions and mathematical statement}
For a group \(G\), let \(\operatorname{Aut}(G)\) be the group of all bijective group homomorphisms from \(G\) to itself. For a cardinal \(\lambda\), write \(2^{\lambda}\) for the cardinality of the power set of a set of cardinality \(\lambda\). A cardinal \(\kappa\) is regular if \(\operatorname{cf}(\kappa)=\kappa\).

The supplied record does not define \(\kappa\)-existential closure. The question below therefore retains that term with its intended source meaning; a complete standalone formulation requires its definition, including the permitted sizes of systems, the treatment of inequations and parameters, and the class of extensions considered. The record also does not explicitly require \(\kappa\) to be infinite.

For every regular cardinal \(\kappa\), every cardinal \(\lambda>\kappa\), and every \(\kappa\)-existentially closed group \(G\) with \(|G|=\lambda\), must one have
\[
\bigl|\operatorname{Aut}(G)\bigr|=2^{\lambda}?
\]

The supplied record reports an affirmative answer in the boundary case \(\lambda=\kappa\), citing the 2025 article listed in Sources. That case lies outside the strict inequality in the question. This reported result and the supplied \emph{open} label are source assertions, not independently verified assessments.
\subsection{Short English statement}
Must every \(\kappa\)-existentially closed group of cardinality \(\lambda>\kappa\), with \(\kappa\) regular, have exactly \(2^{\lambda}\) automorphisms? The source leaves the definition of \(\kappa\)-existential closure unstated and reports that equality is known when \(\lambda=\kappa\).
\subsection{Sources}
\begin{itemize}
\item Local UnsolvedMath registry, permanent ID \(701698\), source code \texttt{catalogue entry}, title \emph{Kourovka Notebook Problem 20.40}. The supplied record labels the problem open and gives the proposed year as 2022.
\item \emph{The Kourovka Notebook (Unsolved Problems in Group Theory)}, Problem 20.40. The supplied background identifies the 20th issue (2022) and the 21st issue edition (2026), \texttt{arXiv:1401.0300v48}. The collection description names E. I. Khukhro and V. D. Mazurov as editors.
\item Proposed attribution as supplied: B. Kaya and a second author whose name is truncated in the attribution field. The background also contains a displaced ending of that name. The article reference below supplies the spelling M. Kuzucuo\u{g}lu; the damaged proposer attribution requires checking.
\item B. Kaya, M. Kuzucuo\u{g}lu, P. Longobardi, M. Maj, \emph{J. Algebra} \textbf{666} (2025), 840--849. The supplied statement cites this article for an affirmative answer when \(\lambda=\kappa\). No article title, DOI, or URL is supplied.
\item These references and the reported literature assessment are reproduced from the supplied record; no independent source consultation or mathematical verification is asserted.
\end{itemize}
\end{document}
